\section{The Saints Who Saw}

The Church has received, among the testimonies of her saints, two accounts of hell that have shaped the prayer of many Christians: the vision of Saint Teresa of Jesus, which she recorded in the thirty-second chapter of her \emph{Life}, and the vision shown to three children at Fatima on 13 July 1917. Both are private revelations. The Church does not propose them for divine faith; she permits the faithful to receive them with prudence, as a help to the faith she teaches.\endnote{On the theological status of private revelation, see Congregation for the Doctrine of the Faith, \emph{The Message of Fatima} (26 June 2000), theological commentary, official Vatican English, vatican.va, accessed 16 October 2026.} Teresa's account is the testimony of a canonized saint and Doctor of the Church; the visions of Fatima were declared worthy of belief by the Bishop of Leiria in 1930, and the Holy See published Sister Lúcia's account of them in 2000.\endnote{Bishop of Leiria, pastoral letter of 13 October 1930, \emph{Documentação Crítica de Fátima: Seleção de documentos (1917--1930)} (Fátima: Santuário de Fátima, 2013), doc.~133; Congregation for the Doctrine of the Faith, \emph{The Message of Fatima}.} What they add is the witness of those who were made to feel it, and the fruit that the feeling bore.

\subsection{Teresa}

Teresa was already advanced in prayer when it happened. ``I was one day in prayer when I found myself in a moment, without knowing how, plunged apparently into hell. I understood that it was our Lord's will I should see the place which the devils kept in readiness for me, and which I had deserved by my sins. It was but a moment, but it seems to me impossible I should ever forget it even if I were to live many years.''\endnote{Teresa of Jesus, \emph{The Life of St.\ Teresa of Jesus}, ch.~32, \S1, trans.\ David Lewis, third edition enlarged (London: Thomas Baker, 1904).} She describes a narrow, dark, foul passage, and a hole in the wall in which she found herself confined. Then she turns from what she saw to what she felt:

\begin{quote}
These sufferings were nothing in comparison with the anguish of my soul, a sense of oppression, of stifling, and of pain so keen, accompanied by so hopeless and cruel an infliction, that I know not how to speak of it. If I said that the soul is continually being torn from the body, it would be nothing, for that implies the destruction of life by the hands of another; but here it is the soul itself that is tearing itself in pieces. I cannot describe that inward fire or that despair, surpassing all torments and all pain\ldots{} There was no light, but all was thick darkness. I do not understand how it is; though there was no light, yet everything that can give pain by being seen was visible.\endnote{Ibid., \S\S4--5.}
\end{quote}

\noindent Everything the Fathers taught is there: the fire within, the despair, the darkness that is not dark enough to hide what torments. And Teresa, who had read of hell in books and meditated on its pains, found her reading surpassed: ``the one is a reality, the other a picture; and all burning here in this life is as nothing in comparison with the fire that is there.''\endnote{Ibid., \S6.}

But the vision did not leave her in terror. She understood it as mercy: ``I understood distinctly that it was a great mercy that our Lord would have me see with mine own eyes the very place from which His compassion saved me.''\endnote{Ibid.} It freed her from every lesser fear, so that ``everything seems endurable in comparison with one instant of suffering such as those I had then to bear in hell.''\endnote{Ibid., \S8.} And it turned her whole heart outward, toward others in danger:

\begin{quote}
It was that vision that filled me with the very great distress which I feel at the sight of so many lost souls\ldots{} and also gave me the most vehement desires for the salvation of souls; for certainly I believe that, to save even one from those overwhelming torments, I would most willingly endure many deaths.\endnote{Ibid., \S9.}
\end{quote}

\noindent She drew the practical conclusion for every Christian: ``Let us, then, for the love of God, avoid all occasions of sin, and our Lord will help us, as He has helped me.''\endnote{Ibid., \S10.} In the same chapter she passes straight from the vision to the longing for penance that it left in her, and from that longing to the foundation of the monastery of Saint Joseph at Avila, the first house of her reform.\endnote{Ibid., \S\S11--13 and the chapter heading.}

\subsection{Fatima}

Sister Lúcia, the surviving witness, wrote down the first part of the secret of 13 July 1917 in her third memoir. ``Our Lady showed us a great sea of fire which seemed to be under the earth,'' she wrote, and she told how the vision filled them with fear. Then the Mother of God spoke to them, ``so kindly and so sadly'':

\begin{quote}
You have seen hell where the souls of poor sinners go. To save them, God wishes to establish in the world devotion to my Immaculate Heart. If what I say to you is done, many souls will be saved and there will be peace.\endnote{Sister Lúcia, Third Memoir (31 August 1941), as published in Congregation for the Doctrine of the Faith, \emph{The Message of Fatima} (2000), official Vatican English, vatican.va, accessed 16 October 2026.}
\end{quote}

\noindent The vision lasted, Lúcia says, but an instant. Our Lady's first words state a fact, and her next a remedy. She showed the children where sinners go in order that they might help them not to go there, and she named the remedy that God wished to give: devotion to her Immaculate Heart. In the third part of the same secret the angel cries out three times, ``Penance, Penance, Penance!''\endnote{Ibid., third part of the secret.}

\subsection{The fruit of seeing}

The two visions, centuries apart, bear the same fruit. Neither leaves the soul staring into the fire. Both turn at once to intercession. Gregory the Great had taught that the saints pray for sinners ``at such time as their hearts may be turned to fruitful penance, and so be saved,''\endnote{Gregory, \emph{Dialogues} IV.44.} and that time is now, while those for whom we pray are still on the way. A Christian who has truly considered hell cannot be indifferent to the souls around him: the colleague who has given up the faith, the child who has stopped going to Mass, the friend who lives in mortal sin and laughs at confession. Teresa would have died many deaths for one of them. The least we can do is pray, offer a sacrifice, and speak a word when charity and prudence allow.
